\pdfbookmark[1]{Abstract}{Abstract}
\chapter*{Abstract}

\textbf{The subject is the building block and the rule.} An atom is a point carrying values: a
count of protons, a count of valence electrons, a radius, an electronegativity. A molecule or a
solid is a set of such points in space. What tells that set from a random packing of the same atoms
is one rule, valence, and valence fixes a bond to a length before anyone measures it. The engine
this book is written for reads any domain as points carrying values and asks one question of the
arrangement: how far it sits from the most disordered arrangement of its own parts. Chemistry is
that question with the atom as the part.

\textbf{What is quoted and what is new.} The building blocks and the rules in the first four
chapters are the settled content of general chemistry and are cited in the last: the periodic law,
the shell counts, the octet, the electronegativity scale, the covalent and van der Waals radii, and
the bond lengths tabulated from crystallography and gas-phase spectroscopy. Nothing about them is
claimed here. What this book adds is their placement in the engine's six parts, and the predictions
that placement makes. The atom is a representation, the octet is a sift, the disordered packing is a
reference, and the bond length is an oracle of the strongest kind, a fact rather than a consensus.

\textbf{Nothing in this book has been measured here.} The engine has readers for a protein and for a
crystal, and both already rest on chemistry: a backbone repeats at three because three bond lengths
are fixed, and a cell edge is a bond geometry tiled. It has no reader for a molecule as such, and
the oracle README reserves the directory chemistry's bond lengths will occupy when they arrive. This
book states what such a reader would find, so the predictions are on the page before the reader
exists. The one demonstration that runs today on the primitives already in the tree is in
\texttt{examples/chemistry}, and it is named there as what it is.

\textbf{Why no set here is bounded.} The engine's own record is that every bound put on a quantity
has had to be taken back off, and the reading improved each time. Chemistry invites the same error
in a sharper form, because a bond either exists or it does not and a tolerance is the obvious way to
decide. This book keeps to the engine's discipline instead: a departure is measured against a
background drawn from the atoms themselves, and no threshold, no cutoff radius and no bounded element
set is chosen by judgment. The question asked is always how far, never whether.
